% 54e Olympiade Internationale de Mathématiques — Santa Marta, Colombie.
% Énoncés officiels en français (imo-official.org). Chaque problème vaut 7 points.

\begin{center}
  \emph{Premier jour --- mardi 23 juillet 2013 --- Durée : 4 heures 30}
\end{center}

\subsection*{Problème 1}
Montrer que pour toute paire d'entiers strictement positifs, $k$ et $n$, il existe $k$ entiers strictement
positifs $m_1, m_2, \ldots, m_k$ (non nécessairement distincts) tels que
\[
1 + \frac{2^k - 1}{n} \ = \ \left(1 + \frac{1}{m_1}\right)\left(1 + \frac{1}{m_2}\right)\cdots\left(1 + \frac{1}{m_k}\right).
\]

\subsection*{Problème 2}
Une configuration de $4027$ points du plan est appelée \emph{colombienne} si elle est constituée de $2013$
points de couleur rouge et de $2014$ points de couleur bleue, et si trois quelconques de ces points ne sont
pas alignés. En traçant des droites, le plan est divisé en régions. Un tracé de droites est appelé \emph{bon}
pour une configuration colombienne si les deux conditions suivantes sont vérifiées :
\begin{itemize}
  \item aucune droite tracée ne passe par un point de la configuration ;
  \item aucune région ne contient des points de couleurs différentes.
\end{itemize}
Trouver la plus petite valeur de $k$ telle que, pour chaque configuration colombienne de $4027$ points, il
existe un bon tracé de $k$ droites.

\subsection*{Problème 3}
Soit $ABC$ un triangle. Le cercle exinscrit au triangle $ABC$ opposé au sommet $A$ est tangent au côté $BC$ au
point $A_1$. Les points $B_1$ sur $CA$ et $C_1$ sur $AB$ sont définis de la même façon, en utilisant les
cercles exinscrits opposés à $B$ et $C$ respectivement. On suppose que le centre du cercle circonscrit au
triangle $A_1B_1C_1$ est un point du cercle circonscrit au triangle $ABC$. Montrer que le triangle $ABC$ est
rectangle.

\emph{Le cercle exinscrit au triangle $ABC$ opposé au sommet $A$ est le cercle tangent au segment $[BC]$, à la
demi-droite $[AB]$ au-delà de $B$ et à la demi-droite $[AC]$ au-delà de $C$. Les cercles exinscrits opposés à
$B$ et à $C$ sont définis de la même manière.}

\bigskip
\begin{center}
  \emph{Second jour --- mercredi 24 juillet 2013 --- Durée : 4 heures 30}
\end{center}

\subsection*{Problème 4}
Soit $ABC$ un triangle d'orthocentre $H$ dont tous les angles sont aigus ; soit $W$ un point du côté $BC$,
compris strictement entre $B$ et $C$. Les points $M$ et $N$ sont, respectivement, les pieds des hauteurs
issues de $B$ et de $C$. On note $\omega_1$ le cercle circonscrit au triangle $BWN$ et $X$ le point de
$\omega_1$ tel que $[WX]$ en soit un diamètre. De la même façon, on note $\omega_2$ le cercle circonscrit au
triangle $CWM$ et $Y$ le point de $\omega_2$ tel que $[WY]$ en soit un diamètre. Montrer que les points $X$,
$Y$ et $H$ sont alignés.

\subsection*{Problème 5}
On note $\mathbb{Q}_{>0}$ l'ensemble des nombres rationnels strictement positifs. Soit
$f : \mathbb{Q}_{>0} \to \mathbb{R}$ une fonction satisfaisant les trois conditions suivantes :
\begin{enumerate}
  \item[(i)] pour tous $x, y \in \mathbb{Q}_{>0}$, on a $f(x)f(y) \geq f(xy)$ ;
  \item[(ii)] pour tous $x, y \in \mathbb{Q}_{>0}$, on a $f(x + y) \geq f(x) + f(y)$ ;
  \item[(iii)] il existe un nombre rationnel $a > 1$ tel que $f(a) = a$.
\end{enumerate}
Montrer que pour tout $x \in \mathbb{Q}_{>0}$, on a $f(x) = x$.

\subsection*{Problème 6}
Soit $n \geq 3$ un entier. On considère un cercle et $n + 1$ points de ce cercle également espacés.

Lors d'un \emph{étiquetage}, chaque point reçoit une étiquette portant l'un des nombres $0, 1, \ldots, n$ de
telle sorte que chaque nombre soit exactement utilisé une fois et une seule. Deux étiquetages sont considérés
comme les mêmes si l'un peut être obtenu à partir de l'autre par une rotation du cercle. Un étiquetage est
appelé \emph{beau} si pour quatre nombres quelconques $a < b < c < d$ figurant sur des étiquettes avec
$a + d = b + c$, la corde joignant les points étiquetés $a$ et $d$ ne coupe pas la corde joignant les points
étiquetés $b$ et $c$.

Soit $M$ le nombre d'étiquetages qui sont beaux, et soit $N$ le nombre de couples $(x, y)$ d'entiers
strictement positifs tels que $x + y \leq n$ et $\mathrm{pgcd}(x, y) = 1$. Montrer que
\[
M = N + 1.
\]
